\subsection{The bounded-test principle}
The finite-radius mechanism in this section is independent of monomials.
We state it before its confluent specialization. This formulation and the
normalization below were isolated in the preceding mathematical audit
\cite{A1v14audit}; we retain that attribution. The result is a bounded
right-inverse principle, not a second positive-history attainment theorem.

\begin{proposition}[Relative prior perturbations under an observation map]
\label{prop:general-bounded-dual}
Let $\mu$ be a probability, $0<c\le\ell\le C<\infty$, and
$\nu=\ell\mu/(\mu\ell)$. Let $\phi=(\phi_1,\ldots,\phi_q)$ be
bounded real tests with $\|\phi_j\|_\infty\le B$ and
$\operatorname{Cov}_\nu(\phi,\phi)\ge\lambda I_q$, $\lambda>0$.
For $0<\varepsilon\le1/2$, vary $\mu'$ over
\[
 \{(1+u)\mu:\mu u=0,\ \|u\|_\infty\le\varepsilon\},
 \qquad \nu'=\ell\mu'/(\mu'\ell).
\]
For every linear map $A:\R^q\to\R^d$ the displacement set satisfies
\[
 c_0\varepsilon A[-1,1]^q
 \ \subset\ \{A(\nu'\phi-\nu\phi)\}
 \ \subset\ C_0\varepsilon A[-1,1]^q,
\]
where $c_0,C_0>0$ depend only on $q,B,\lambda$, not on the singular
values or rank of $A$.
\end{proposition}
\begin{proof}
For $\mu'=(1+u)\mu$, the exact identity
\[
 \nu'\phi-\nu\phi
 =\frac{\nu\{u(\phi-\nu\phi)\}}{1+\nu u}
\]
bounds each coordinate by $4B\varepsilon$. For the other inclusion put
$\Gamma=\operatorname{Cov}_\nu(\phi,\phi)$ and, for
$v\in[-1,1]^q$, set
\[
 f_v=v^{\mathsf T}\Gamma^{-1}(\phi-\nu\phi),\qquad
 K=\max\{1,2qB/\lambda\},\qquad s=\varepsilon/(4K).
\]
Then $\nu f_v=0$, $\nu(\phi f_v)=v$, and $\|f_v\|_\infty\le K$.
Define a prior, rather than only a posterior perturbation, by
\[
 \frac{d\mu'}{d\mu}=\frac{1+s f_v}{1+s\mu f_v}.
\]
It has total mass one, is positive, and obeys
\[
 \left\|\frac{d\mu'}{d\mu}-1\right\|_\infty
 \le\frac{2sK}{1-sK}
 =\frac{\varepsilon/2}{1-\varepsilon/4}\le\varepsilon.
\]
The evidence is $\mu\ell/(1+s\mu f_v)$, so Bayes normalization gives
exactly $\nu'=(1+s f_v)\nu$ and hence displacement $sv$. This realizes
one common cube, not merely separate directional alternatives. Applying
$A$ to both inclusions completes the proof, also when $A$ has a kernel.
\end{proof}

\subsection{Complete confluent flags and all width orders}
